\documentclass[aspectratio=169,10pt]{beamer}

% ============================================================
% UNIT-IV: WORK STUDY AND STATISTICAL QUALITY CONTROL
% FULL DETAILED METROPOLIS BEAMER PRESENTATION
% PDFLaTeX / OVERLEAF FREE COMPATIBLE
% ============================================================

\usetheme[
    progressbar=frametitle,
    numbering=fraction
]{metropolis}

% ------------------------------------------------------------
% PDFLaTeX-SAFE FONT CONFIGURATION
% ------------------------------------------------------------
\usepackage[T1]{fontenc}
\usepackage[utf8]{inputenc}
\usepackage{lmodern}

% ------------------------------------------------------------
% MATHEMATICS
% ------------------------------------------------------------
\usepackage{amsmath}
\usepackage{amssymb}
\usepackage{mathtools}

% ------------------------------------------------------------
% TABLES
% ------------------------------------------------------------
\usepackage{booktabs}
\usepackage{array}
\usepackage{tabularx}
\usepackage{multirow}

% ------------------------------------------------------------
% GRAPHICS / FORMATTING
% ------------------------------------------------------------
\usepackage{graphicx}
\usepackage{microtype}
\usepackage{enumitem}

% ------------------------------------------------------------
% COLORS
% ------------------------------------------------------------
\definecolor{HeaderDark}{RGB}{35,39,42}
\definecolor{Accent}{RGB}{0,150,136}
\definecolor{TextDark}{RGB}{45,45,45}
\definecolor{Muted}{RGB}{100,105,110}
\definecolor{LightAccent}{RGB}{232,245,243}
\definecolor{DarkBlue}{RGB}{35,70,105}
\definecolor{DarkGreen}{RGB}{30,105,75}
\definecolor{DarkRed}{RGB}{145,45,50}

\setbeamercolor{normal text}{fg=TextDark}
\setbeamercolor{frametitle}{fg=HeaderDark}
\setbeamercolor{progress bar}{fg=Accent}
\setbeamercolor{alerted text}{fg=Accent}
\setbeamercolor{example text}{fg=DarkGreen}

% ------------------------------------------------------------
% FONT SIZES
% ------------------------------------------------------------
\setbeamerfont{frametitle}{size=\large}
\setbeamerfont{framesubtitle}{size=\normalsize}
\setbeamerfont{title}{size=\LARGE}
\setbeamerfont{subtitle}{size=\large}

% ------------------------------------------------------------
% LIST SPACING
% ------------------------------------------------------------
\setlist[itemize]{
    topsep=2pt,
    itemsep=2pt,
    parsep=0pt,
    partopsep=0pt
}

\setlist[enumerate]{
    topsep=2pt,
    itemsep=2pt,
    parsep=0pt,
    partopsep=0pt
}

% ------------------------------------------------------------
% MATHEMATICAL SHORTCUTS
% ------------------------------------------------------------
\newcommand{\OT}{T_o}
\newcommand{\RT}{R}
\newcommand{\NT}{T_N}
\newcommand{\ST}{T_S}
\newcommand{\Allow}{A}
\newcommand{\xbar}{\bar{x}}
\newcommand{\Rbar}{\bar{R}}
\newcommand{\phat}{\hat{p}}
\newcommand{\qhat}{\hat{q}}
\newcommand{\Pa}{P_a}
\newcommand{\LTPD}{\mathrm{LTPD}}
\newcommand{\AQL}{\mathrm{AQL}}

% ------------------------------------------------------------
% TABLE COLUMN TYPES
% ------------------------------------------------------------
\newcolumntype{Y}{>{\raggedright\arraybackslash}X}
\newcolumntype{C}{>{\centering\arraybackslash}X}

% ------------------------------------------------------------
% TITLE INFORMATION
% ------------------------------------------------------------
\title{UNIT--IV: WORK STUDY AND\\
STATISTICAL QUALITY CONTROL}

\subtitle{
Detailed Theory, Mathematical Modelling,\\
Numerical Examples, Interpretation and Applications
}

\author{Gajavada Sanjeevkumar}

\institute{
Department of Mechanical Engineering\\
Indian Institute of Technology (Indian School of Mines), Dhanbad
}

\date{}

% ============================================================
% DOCUMENT
% ============================================================

\begin{document}

% ============================================================
% TITLE
% ============================================================

\maketitle

% ============================================================
% CONTENTS
% ============================================================

\begin{frame}{Unit--IV Roadmap}
\small
\begin{enumerate}
    \item Work Study
    \item Method Study
    \item Method Study Charts
    \item Work Measurement
    \item Stopwatch Time Study
    \item Performance Rating and Normal Time
    \item Allowances and Standard Time
    \item Work Sampling
    \item Statistical Quality Control
    \item Control Charts
    \item Variables and Attributes Charts
    \item $\bar X$--$R$ Charts
    \item $p$ and $c$ Charts
    \item Acceptance Sampling
    \item Operating Characteristic Curves
    \item AQL, LTPD and Sampling Risks
    \item Integrated Manufacturing Example
    \item Formula Sheet
    \item Viva Questions
\end{enumerate}
\end{frame}

% ============================================================
% SECTION I
% ============================================================

\section{Work Study}

\begin{frame}{What is Work Study?}
\small

\textbf{Work study} is a systematic examination of the methods of doing work
and the time required to perform the work.

\vspace{4pt}

It consists primarily of:

\begin{itemize}
    \item \textbf{Method study} -- determines the best practical way of doing a job.
    \item \textbf{Work measurement} -- determines the time required for a qualified worker
    to perform the job at a defined level of performance.
\end{itemize}

\vspace{4pt}

Therefore,

\[
\boxed{
\text{Work Study}
=
\text{Method Study}
+
\text{Work Measurement}
}
\]

\textbf{Central objective:}

\[
\boxed{
\text{Improve productivity by eliminating unnecessary work and establishing standards}
}
\]

\end{frame}

% ------------------------------------------------------------

\begin{frame}{Why is Work Study Required?}
\small

Work study is required because manufacturing systems contain:

\begin{itemize}
    \item unnecessary motions,
    \item unnecessary transportation,
    \item waiting and delays,
    \item poor workplace layout,
    \item unbalanced operations,
    \item inefficient material handling,
    \item inappropriate working methods,
    \item excessive fatigue,
    \item uncontrolled variability in task times.
\end{itemize}

\vspace{4pt}

The fundamental engineering question is:

\begin{center}
\Large
\textbf{Can the same useful output be obtained with less unnecessary effort, time and cost?}
\end{center}

\end{frame}

% ------------------------------------------------------------

\begin{frame}{Objectives of Work Study}
\small

\begin{columns}[T]

\column{0.48\textwidth}

\textbf{Technical objectives}

\begin{itemize}
    \item Improve methods.
    \item Eliminate unnecessary operations.
    \item Reduce material movement.
    \item Improve workplace layout.
    \item Reduce cycle time.
    \item Establish standard methods.
\end{itemize}

\column{0.48\textwidth}

\textbf{Managerial objectives}

\begin{itemize}
    \item Production planning.
    \item Labour planning.
    \item Capacity planning.
    \item Cost estimation.
    \item Incentive schemes.
    \item Productivity improvement.
\end{itemize}

\end{columns}

\vspace{6pt}

\[
\boxed{
\text{Better method}
\rightarrow
\text{Less time}
\rightarrow
\text{Higher productivity}
\rightarrow
\text{Lower unit cost}
}
\]

\end{frame}

% ------------------------------------------------------------

\begin{frame}{Work Study: Basic Procedure}
\small

A systematic work-study investigation generally follows:

\begin{enumerate}
    \item Select the job.
    \item Record the existing method.
    \item Examine the recorded facts.
    \item Develop an improved method.
    \item Install the improved method.
    \item Measure the work.
    \item Determine the standard time.
    \item Maintain the improved standard.
\end{enumerate}

\vspace{5pt}

\textbf{Important principle:}

The analyst should not immediately measure the old method if the method
itself can first be improved.

\[
\boxed{
\text{Improve the method first, then establish the time standard}
}
\]

\end{frame}

% ============================================================
% METHOD STUDY
% ============================================================

\section{Method Study}

\begin{frame}{Method Study}
\small

\textbf{Method study} is the systematic recording and critical examination
of existing and proposed ways of doing work.

\vspace{5pt}

It attempts to answer:

\begin{itemize}
    \item What operation is being performed?
    \item Why is it necessary?
    \item Where is it performed?
    \item Why is it performed there?
    \item When is it performed?
    \item Who performs it?
    \item Why is that person doing it?
    \item How is it performed?
\end{itemize}

\vspace{5pt}

The aim is not simply to make a worker work faster.

\[
\boxed{
\text{Aim = eliminate unnecessary work and improve the system}
}
\]

\end{frame}

% ------------------------------------------------------------

\begin{frame}{Method Study Procedure}
\small

A commonly used sequence is:

\begin{enumerate}
    \item \textbf{Select} the work to be studied.
    \item \textbf{Record} all relevant facts.
    \item \textbf{Examine} the facts critically.
    \item \textbf{Develop} an improved method.
    \item \textbf{Evaluate} the proposed method.
    \item \textbf{Define} the new method.
    \item \textbf{Install} the method.
    \item \textbf{Maintain} the method.
\end{enumerate}

\vspace{5pt}

The examination stage is especially important because observations
without critical questioning do not automatically produce improvement.

\end{frame}

% ------------------------------------------------------------

\begin{frame}{SREDIM Technique}
\small

\textbf{SREDIM} is a traditional mnemonic for method-study procedure:

\[
\boxed{
S-R-E-D-I-M
}
\]

\begin{description}
    \item[S] Select
    \item[R] Record
    \item[E] Examine
    \item[D] Develop
    \item[I] Install
    \item[M] Maintain
\end{description}

\vspace{4pt}

\textbf{Interpretation:}

\begin{itemize}
    \item \textbf{Select}: identify the job.
    \item \textbf{Record}: document the existing method.
    \item \textbf{Examine}: question each activity.
    \item \textbf{Develop}: formulate alternatives.
    \item \textbf{Install}: introduce the selected method.
    \item \textbf{Maintain}: ensure continued adherence and improvement.
\end{itemize}

\end{frame}

% ------------------------------------------------------------

\begin{frame}{Critical Examination Questions}
\small

For each activity, ask:

\begin{center}
\Large
\textbf{WHY?}
\end{center}

Then investigate:

\begin{tabularx}{\textwidth}{YY}
\toprule
Question & Engineering purpose \\
\midrule
What? & Identify the exact activity.\\
Why? & Test whether it is necessary.\\
Where? & Check location and layout.\\
When? & Check sequence and timing.\\
Who? & Check responsibility and skill.\\
How? & Check the physical method.\\
\bottomrule
\end{tabularx}

\vspace{4pt}

The goal is to distinguish \textbf{necessary work} from
\textbf{avoidable work}.

\end{frame}

% ------------------------------------------------------------

\begin{frame}{ECRS Principle}
\small

ECRS provides four basic directions for method improvement.

\[
\boxed{
E-C-R-S
}
\]

\begin{description}
    \item[E -- Eliminate] Remove unnecessary activities.
    \item[C -- Combine] Combine activities where practical.
    \item[R -- Rearrange] Change sequence or arrangement.
    \item[S -- Simplify] Make the remaining work easier.
\end{description}

\vspace{5pt}

\textbf{Example:}

Suppose a worker:

\[
\text{Walk} \rightarrow \text{Pick} \rightarrow
\text{Walk} \rightarrow \text{Place}
\]

A better layout may permit:

\[
\text{Pick} \rightarrow \text{Place}
\]

with reduced travel.

\end{frame}

% ============================================================
% METHOD STUDY CHARTS
% ============================================================

\section{Method Study Charts}

\begin{frame}{Purpose of Method Study Charts}
\small

Charts convert a physical work process into a structured representation.

They help identify:

\begin{itemize}
    \item operations,
    \item inspections,
    \item transport,
    \item delays,
    \item storage,
    \item unnecessary movement,
    \item repeated activities,
    \item excessive travel,
    \item poor sequence.
\end{itemize}

\vspace{5pt}

Different charts provide different levels of detail.

\[
\boxed{
\text{Overall process}
\rightarrow
\text{Detailed process}
\rightarrow
\text{Operator-machine interaction}
}
\]

\end{frame}

% ------------------------------------------------------------

\begin{frame}{Operation Process Chart}
\small

An \textbf{Operation Process Chart (OPC)} records the principal
operations and inspections involved in producing a product.

\vspace{5pt}

Typical symbols:

\begin{itemize}
    \item Operation: circle
    \item Inspection: square
\end{itemize}

It is useful for:

\begin{itemize}
    \item understanding overall manufacturing sequence,
    \item identifying major operations,
    \item identifying inspection stages,
    \item comparing alternative processes.
\end{itemize}

\vspace{5pt}

\textbf{Limitation:}

It does not provide the same detailed information about transportation,
delays and storage as a flow process chart.

\end{frame}

% ------------------------------------------------------------

\begin{frame}{Flow Process Chart}
\small

A \textbf{Flow Process Chart} gives a more detailed representation.

It normally identifies:

\begin{itemize}
    \item Operation
    \item Inspection
    \item Transport
    \item Delay
    \item Storage
\end{itemize}

\vspace{4pt}

A typical analysis asks:

\[
\text{How many operations?}
\]

\[
\text{How much distance is travelled?}
\]

\[
\text{How much waiting occurs?}
\]

\[
\text{How many storage points exist?}
\]

\vspace{4pt}

It is particularly useful for material handling and layout improvement.

\end{frame}

% ------------------------------------------------------------

\begin{frame}{Two-Handed Process Chart}
\small

The two-handed process chart records the activities of:

\[
\boxed{\text{Left hand and right hand}}
\]

during a short repetitive task.

\vspace{5pt}

It is useful for:

\begin{itemize}
    \item repetitive assembly,
    \item manual operations,
    \item identifying idle hand time,
    \item balancing hand motions,
    \item improving workplace arrangement.
\end{itemize}

\vspace{5pt}

\textbf{Ideal principle:}

Both hands should have purposeful activity where practical,
while avoiding unnecessary motion.

\end{frame}

% ------------------------------------------------------------

\begin{frame}{Man--Machine Chart}
\small

A \textbf{Man--Machine Chart} represents the utilization of a worker
and one or more machines over a common time scale.

\vspace{5pt}

It identifies:

\begin{itemize}
    \item worker working time,
    \item worker idle time,
    \item machine running time,
    \item machine idle time,
    \item simultaneous utilization.
\end{itemize}

\vspace{5pt}

A major objective is to reduce idle resources.

For example:

\[
\boxed{
\text{Worker idle while machine runs}
}
\]

may provide an opportunity for another task.

\end{frame}

% ------------------------------------------------------------

\begin{frame}{Comparison of Method Study Charts}
\small

\begin{tabularx}{\textwidth}{YYYY}
\toprule
Chart & Main purpose & Detail & Typical use \\
\midrule
OPC & Major operations & Low--medium & Overall process\\
Flow Process Chart & Detailed flow & High & Material/process analysis\\
Two-Handed Chart & Hand motions & Very high & Manual repetitive work\\
Man--Machine Chart & Resource utilization & High & Worker-machine analysis\\
\bottomrule
\end{tabularx}

\vspace{8pt}

\textbf{Selection rule:}

Choose the chart according to the question being investigated.

\end{frame}

% ============================================================
% WORK MEASUREMENT
% ============================================================

\section{Work Measurement}

\begin{frame}{What is Work Measurement?}
\small

\textbf{Work measurement} is the application of techniques designed
to establish the time required by a qualified worker to complete
a specified task at a defined level of performance.

\vspace{5pt}

The result is generally a:

\[
\boxed{\text{Standard Time}}
\]

Standard time is useful for:

\begin{itemize}
    \item production planning,
    \item scheduling,
    \item manpower estimation,
    \item costing,
    \item capacity planning,
    \item incentive systems,
    \item productivity measurement.
\end{itemize}

\end{frame}

% ------------------------------------------------------------

\begin{frame}{Major Work Measurement Techniques}
\small

Important techniques include:

\begin{enumerate}
    \item Stopwatch time study
    \item Work sampling
    \item Predetermined motion-time systems
    \item Standard data systems
    \item Analytical estimating
\end{enumerate}

\vspace{6pt}

For repetitive manufacturing work, stopwatch time study is widely used.

For work involving a large number of random activities,
work sampling can be more practical.

\end{frame}

% ------------------------------------------------------------

\begin{frame}{Time Study: Basic Concept}
\small

Time study determines how long a task should take under specified
conditions.

The analyst:

\begin{enumerate}
    \item Selects the task.
    \item Divides it into elements.
    \item Observes several cycles.
    \item Records element times.
    \item Determines representative observed time.
    \item Applies performance rating.
    \item Adds appropriate allowances.
\end{enumerate}

The final objective is:

\[
\boxed{
\text{Observed Time}
\rightarrow
\text{Normal Time}
\rightarrow
\text{Standard Time}
}
\]

\end{frame}

% ============================================================
% STOPWATCH TIME STUDY
% ============================================================

\section{Stopwatch Time Study}

\begin{frame}{Stopwatch Time Study}
\small

Stopwatch time study is a direct work-measurement technique.

The analyst observes a worker performing a defined task and measures
the time associated with individual elements or complete cycles.

\vspace{5pt}

\textbf{Suitable when:}

\begin{itemize}
    \item the work is repetitive,
    \item cycle boundaries are identifiable,
    \item the method is relatively stable,
    \item sufficient observations can be obtained.
\end{itemize}

\end{frame}

% ------------------------------------------------------------

\begin{frame}{Why Divide a Job into Elements?}
\small

A complete job may contain several distinct activities.

Example:

\[
\text{Pick component}
\rightarrow
\text{Position}
\rightarrow
\text{Fasten}
\rightarrow
\text{Inspect}
\]

Dividing the job provides:

\begin{itemize}
    \item better observation accuracy,
    \item identification of abnormal elements,
    \item element-wise rating,
    \item identification of bottlenecks,
    \item better method improvement.
\end{itemize}

\vspace{5pt}

\[
\boxed{
T_{\text{cycle}}
=
\sum_{i=1}^{m} T_i
}
\]

\end{frame}

% ------------------------------------------------------------

\begin{frame}{Continuous Timing}
\small

In \textbf{continuous timing}, the stopwatch runs continuously.

The observer records the stopwatch reading at the end of each element.

If cumulative readings are:

\[
C_1,\ C_2,\ C_3,\ldots,C_n
\]

then the element times are:

\[
T_1=C_1
\]

\[
T_2=C_2-C_1
\]

\[
T_3=C_3-C_2
\]

and generally:

\[
\boxed{
T_i=C_i-C_{i-1}
}
\]

with \(C_0=0\).

\end{frame}

% ------------------------------------------------------------

\begin{frame}{Snapback / Flyback Timing}
\small

In snapback timing, the stopwatch is reset after each element.

Therefore the displayed reading directly represents the element time.

Advantages:

\begin{itemize}
    \item easy element-wise reading,
    \item direct recording,
    \item convenient for repetitive cycles.
\end{itemize}

Limitations:

\begin{itemize}
    \item resetting introduces observer reaction effects,
    \item small timing errors may accumulate,
    \item requires careful observer technique.
\end{itemize}

\end{frame}

% ------------------------------------------------------------

\begin{frame}{Observed Time}
\small

Suppose an element is observed for \(n\) cycles:

\[
T_1,T_2,\ldots,T_n
\]

The average observed time is:

\[
\boxed{
\bar T_o
=
\frac{1}{n}
\sum_{i=1}^{n}T_i
}
\]

If the observations are:

\[
0.42,\quad0.44,\quad0.43,\quad0.45,\quad0.41
\]

then:

\[
\bar T_o
=
\frac{0.42+0.44+0.43+0.45+0.41}{5}
\]

\[
\boxed{
\bar T_o=0.430\ \text{min}
}
\]

\end{frame}

% ------------------------------------------------------------

\begin{frame}{Abnormal Observations}
\small

Not every observed value should automatically be treated as normal.

Possible causes of abnormal observations:

\begin{itemize}
    \item machine malfunction,
    \item material shortage,
    \item interruption,
    \item unusual fatigue,
    \item defective component,
    \item communication interruption,
    \item measurement error.
\end{itemize}

\vspace{5pt}

The analyst should investigate the reason before deciding whether
an observation should be excluded.

\[
\boxed{
\text{Statistical outlier} \neq \text{automatically invalid observation}
}
\]

\end{frame}

% ============================================================
% PERFORMANCE RATING
% ============================================================

\section{Performance Rating}

\begin{frame}{Why is Performance Rating Needed?}
\small

Two workers may perform the same task at different observed speeds.

Observed time alone therefore does not necessarily represent
the time corresponding to a defined standard performance level.

Performance rating adjusts the observed time.

\[
\boxed{
\text{Normal Time}
=
\text{Observed Time}
\times
\text{Rating Factor}
}
\]

If rating is expressed as a percentage:

\[
\boxed{
NT
=
OT\left(\frac{R}{100}\right)
}
\]

\end{frame}

% ------------------------------------------------------------

\begin{frame}{Physical Meaning of Rating}
\small

Consider:

\[
OT=0.50\text{ min}
\]

and

\[
R=110\%.
\]

Then:

\[
NT
=
0.50
\left(
\frac{110}{100}
\right)
\]

\[
\boxed{
NT=0.55\text{ min}
}
\]

Interpretation:

The worker was observed performing at approximately \(110\%\) of
the defined standard pace.

Therefore the equivalent time at standard performance is greater
than the observed time.

\end{frame}

% ------------------------------------------------------------

\begin{frame}{Rating Less Than 100\%}
\small

Suppose:

\[
OT=0.80\text{ min}
\]

and

\[
R=90\%.
\]

Then:

\[
NT
=
0.80
\left(
\frac{90}{100}
\right)
\]

\[
\boxed{
NT=0.72\text{ min}
}
\]

Because the observed worker operated below the reference performance
level, the equivalent standard-performance time is lower than the
observed time.

\end{frame}

% ------------------------------------------------------------

\begin{frame}{Rating Methods}
\small

Common approaches include:

\begin{itemize}
    \item speed rating,
    \item Westinghouse system,
    \item synthetic rating,
    \item objective rating,
    \item comparison with defined standard performance.
\end{itemize}

A rating system should attempt to improve consistency between observers.

\vspace{5pt}

However, rating contains human judgment.

Therefore:

\[
\boxed{
\text{Rating is a major source of measurement variability}
}
\]

\end{frame}

% ============================================================
% NORMAL TIME
% ============================================================

\section{Normal Time}

\begin{frame}{Normal Time}
\small

The \textbf{normal time} is the time required to perform the task
at the defined standard performance level, excluding allowances.

\[
\boxed{
NT=OT\left(\frac{R}{100}\right)
}
\]

where:

\begin{itemize}
    \item \(OT\) = observed time,
    \item \(R\) = performance rating in percent,
    \item \(NT\) = normal time.
\end{itemize}

\vspace{5pt}

The distinction is:

\[
\boxed{
OT \rightarrow \text{what was observed}
}
\]

\[
\boxed{
NT \rightarrow \text{time normalized to reference performance}
}
\]

\end{frame}

% ------------------------------------------------------------

\begin{frame}{Numerical Example: Normal Time}
\small

A worker completes an element in an average observed time of:

\[
OT=1.20\text{ min}
\]

The performance rating is:

\[
R=95\%.
\]

Therefore:

\[
NT
=
1.20
\left(
\frac{95}{100}
\right)
\]

\[
NT=1.14\text{ min}
\]

Hence:

\[
\boxed{
NT=1.14\text{ min}
}
\]

This is the time corresponding to the reference performance level.

\end{frame}

% ============================================================
% ALLOWANCES
% ============================================================

\section{Allowances and Standard Time}

\begin{frame}{Why are Allowances Required?}
\small

A worker cannot continuously perform productive work for the entire
available working time.

Allowances account for legitimate time associated with:

\begin{itemize}
    \item personal needs,
    \item fatigue,
    \item unavoidable delays,
    \item machine interruptions,
    \item process delays,
    \item special conditions.
\end{itemize}

Therefore:

\[
\boxed{
ST>NT
}
\]

for a positive allowance.

\end{frame}

% ------------------------------------------------------------

\begin{frame}{Types of Allowances}
\small

Important categories include:

\begin{enumerate}
    \item \textbf{Personal allowance}
    \item \textbf{Fatigue allowance}
    \item \textbf{Delay allowance}
    \item \textbf{Contingency allowance}
    \item \textbf{Special allowance}
\end{enumerate}

\vspace{5pt}

The exact allowance depends on:

\begin{itemize}
    \item work environment,
    \item physical effort,
    \item mental effort,
    \item posture,
    \item temperature,
    \item vibration,
    \item lighting,
    \item process characteristics.
\end{itemize}

\end{frame}

% ------------------------------------------------------------

\begin{frame}{Standard Time: Allowance Convention 1}
\small

If allowance is expressed as a fraction of normal time:

\[
A=\frac{\text{Allowance time}}{\text{Normal time}}
\]

then:

\[
\boxed{
ST=NT(1+A)
}
\]

Example:

\[
NT=10\text{ min},\qquad A=15\%=0.15
\]

Therefore:

\[
ST=10(1.15)
\]

\[
\boxed{
ST=11.5\text{ min}
}
\]

\end{frame}

% ------------------------------------------------------------

\begin{frame}{Standard Time: Allowance Convention 2}
\small

Sometimes the allowance is defined as a fraction of the
\textbf{total standard time}.

If:

\[
A=
\frac{\text{Allowance time}}{\text{Standard time}}
\]

then:

\[
NT=ST(1-A)
\]

Therefore:

\[
\boxed{
ST=\frac{NT}{1-A}
}
\]

Example:

\[
NT=10,\qquad A=0.15
\]

\[
ST=\frac{10}{0.85}
\]

\[
\boxed{
ST=11.765\text{ min}
}
\]

\end{frame}

% ------------------------------------------------------------

\begin{frame}{Why Two Allowance Formulas Exist}
\small

The two equations:

\[
ST=NT(1+A)
\]

and

\[
ST=\frac{NT}{1-A}
\]

are both valid under different definitions of allowance.

\vspace{6pt}

Therefore the analyst must always ask:

\[
\boxed{
\text{Allowance is a percentage of WHAT?}
}
\]

\vspace{5pt}

Never apply an allowance formula without checking its definition.

\end{frame}

% ------------------------------------------------------------

\begin{frame}{Complete Time Study Calculation}
\small

Suppose:

\[
OT=2.0\text{ min}
\]

\[
R=110\%
\]

\[
A=15\%
\]

Using the first convention:

\[
NT
=
2.0(1.10)
=
2.20\text{ min}
\]

Then:

\[
ST
=
2.20(1.15)
\]

\[
\boxed{
ST=2.53\text{ min}
}
\]

Thus:

\[
\boxed{
OT\rightarrow NT\rightarrow ST
}
\]

\[
2.00\rightarrow2.20\rightarrow2.53\text{ min}
\]

\end{frame}

% ============================================================
% WORK SAMPLING
% ============================================================

\section{Work Sampling}

\begin{frame}{What is Work Sampling?}
\small

\textbf{Work sampling} estimates the proportion of time spent on
different activities by taking observations at randomly selected
times.

Unlike stopwatch time study, it does not continuously observe the
worker.

\vspace{5pt}

It is based on probability.

For an activity observed \(x\) times in \(n\) observations:

\[
\boxed{
\hat p=\frac{x}{n}
}
\]

where \(\hat p\) estimates the true proportion \(p\).

\end{frame}

% ------------------------------------------------------------

\begin{frame}{When is Work Sampling Useful?}
\small

Work sampling is especially useful when:

\begin{itemize}
    \item work cycles are long,
    \item activities are irregular,
    \item workers perform several different tasks,
    \item continuous timing is expensive,
    \item machine utilization is being studied,
    \item delay proportions are required.
\end{itemize}

Examples:

\begin{itemize}
    \item maintenance personnel,
    \item supervisors,
    \item machine operators,
    \item material handlers,
    \item office activities.
\end{itemize}

\end{frame}

% ------------------------------------------------------------

\begin{frame}{Probability Model for Work Sampling}
\small

Each observation can be considered as a Bernoulli trial:

\[
X=
\begin{cases}
1,&\text{activity is observed}\\
0,&\text{activity is not observed}
\end{cases}
\]

with:

\[
P(X=1)=p
\]

and

\[
P(X=0)=1-p=q.
\]

For \(n\) independent observations:

\[
X\sim\mathrm{Binomial}(n,p).
\]

Therefore:

\[
E(X)=np
\]

and

\[
Var(X)=np(1-p).
\]

\end{frame}

% ------------------------------------------------------------

\begin{frame}{Estimator of Activity Proportion}
\small

If \(x\) observations correspond to the activity:

\[
\boxed{
\hat p=\frac{x}{n}
}
\]

The complementary proportion is:

\[
\boxed{
\hat q=1-\hat p
}
\]

For example:

\[
x=180,\qquad n=1000
\]

then:

\[
\hat p=\frac{180}{1000}=0.18.
\]

Therefore approximately:

\[
\boxed{
18\%
}
\]

of observed time corresponds to the activity.

\end{frame}

% ------------------------------------------------------------

\begin{frame}{Variance and Standard Error}
\small

For the sample proportion:

\[
Var(\hat p)
=
\frac{p(1-p)}{n}
\]

Therefore:

\[
\boxed{
\sigma_{\hat p}
=
\sqrt{
\frac{p(1-p)}{n}
}
}
\]

Since \(p\) is usually unknown, estimate it using \(\hat p\):

\[
\boxed{
SE(\hat p)
=
\sqrt{
\frac{\hat p(1-\hat p)}{n}
}
}
\]

Thus increasing \(n\) reduces uncertainty.

\end{frame}

% ------------------------------------------------------------

\begin{frame}{Confidence Interval for Work Sampling}
\small

For sufficiently large sample size, an approximate confidence interval is:

\[
\boxed{
\hat p
\pm
z_{\alpha/2}
\sqrt{
\frac{\hat p(1-\hat p)}{n}
}
}
\]

For a 95\% confidence interval:

\[
z_{0.025}=1.96
\]

Hence:

\[
\boxed{
CI_{95\%}
=
\hat p
\pm
1.96
\sqrt{
\frac{\hat p(1-\hat p)}{n}
}
}
\]

\end{frame}

% ------------------------------------------------------------

\begin{frame}{Sample Size for Work Sampling}
\small

If the desired half-width of error is \(e\):

\[
e=
z_{\alpha/2}
\sqrt{
\frac{pq}{n}
}
\]

Squaring:

\[
e^2=
z_{\alpha/2}^2
\frac{pq}{n}
\]

Therefore:

\[
\boxed{
n=
\frac{
z_{\alpha/2}^2pq
}{
e^2
}
}
\]

If \(p\) is unknown, the conservative value:

\[
p=q=0.5
\]

is often used because:

\[
pq=0.25
\]

is maximum at \(p=0.5\).

\end{frame}

% ------------------------------------------------------------

\begin{frame}{Numerical Example: Sample Size}
\small

Find \(n\) for:

\[
95\%\text{ confidence}
\]

\[
e=0.05
\]

and assume:

\[
p=0.50.
\]

Then:

\[
n=
\frac{(1.96)^2(0.5)(0.5)}
{(0.05)^2}
\]

\[
n=
\frac{3.8416(0.25)}
{0.0025}
\]

\[
\boxed{
n\approx384.16
}
\]

Therefore:

\[
\boxed{
n=385\text{ observations}
}
\]

after rounding upward.

\end{frame}

% ------------------------------------------------------------

\begin{frame}{Work Sampling: Numerical Example}
\small

Suppose:

\[
n=500
\]

random observations are taken and the worker is found to be idle in:

\[
x=75
\]

observations.

Then:

\[
\hat p=\frac{75}{500}=0.15.
\]

Therefore:

\[
\boxed{
15\%
}
\]

of observations indicate idle time.

For 95\% confidence:

\[
SE=
\sqrt{
\frac{0.15(0.85)}{500}
}
\]

\[
SE\approx0.01597.
\]

\end{frame}

% ------------------------------------------------------------

\begin{frame}{Interpretation of Work Sampling Result}
\small

For the previous example:

\[
CI
=
0.15
\pm
1.96(0.01597)
\]

\[
CI
\approx
0.15\pm0.0313.
\]

Therefore:

\[
\boxed{
0.1187\le p\le0.1813
}
\]

approximately.

Interpretation:

At the stated confidence level and under the sampling assumptions,
the true idle proportion is estimated within this interval.

\end{frame}

% ============================================================
% STATISTICAL QUALITY CONTROL
% ============================================================

\section{Statistical Quality Control}

\begin{frame}{What is Statistical Quality Control?}
\small

\textbf{Statistical Quality Control (SQC)} uses statistical methods
to monitor, control and improve process quality.

The fundamental idea is:

\[
\boxed{
\text{Variation exists in every process}
}
\]

The engineering objective is to distinguish:

\[
\boxed{
\text{Common causes}
}
\]

from:

\[
\boxed{
\text{Special causes}
}
\]

\end{frame}

% ------------------------------------------------------------

\begin{frame}{Why Does Variation Occur?}
\small

Sources of variation may include:

\begin{itemize}
    \item raw material variation,
    \item machine variation,
    \item tool wear,
    \item temperature,
    \item operator differences,
    \item measurement variation,
    \item environmental conditions,
    \item process disturbances.
\end{itemize}

Variation is therefore not automatically evidence of a defective
process.

The key question is:

\[
\boxed{
\text{Is the observed variation consistent with the established process system?}
}
\]

\end{frame}

% ------------------------------------------------------------

\begin{frame}{Common Causes}
\small

\textbf{Common causes} are sources of variation inherent in the
normal operation of the process.

Examples:

\begin{itemize}
    \item normal machine vibration,
    \item small raw-material differences,
    \item ordinary measurement variation,
    \item normal environmental variation.
\end{itemize}

If only common causes are present, the process is statistically
stable around its current operating condition.

\vspace{5pt}

A change in common-cause variation generally requires a
\textbf{system-level change}.

\end{frame}

% ------------------------------------------------------------

\begin{frame}{Special Causes}
\small

\textbf{Special causes} are unusual sources of variation that are
not part of the normal process system.

Examples:

\begin{itemize}
    \item broken tool,
    \item incorrect machine setting,
    \item defective raw material batch,
    \item sensor malfunction,
    \item sudden temperature excursion,
    \item operator or machine disturbance.
\end{itemize}

The objective of control-chart analysis is to identify evidence
that such causes may be present.

\end{frame}

% ------------------------------------------------------------

\begin{frame}{Control Limits vs Specification Limits}
\small

These concepts must not be confused.

\textbf{Control limits} are calculated from process data.

They describe expected process variation under statistical control.

\textbf{Specification limits} are engineering/customer requirements.

\[
\boxed{
\text{Control limits}\neq\text{Specification limits}
}
\]

A process can be:

\begin{itemize}
    \item statistically stable but incapable of meeting specifications,
    \item unstable but temporarily producing conforming items,
    \item stable and capable,
    \item unstable and incapable.
\end{itemize}

\end{frame}

% ============================================================
% VARIABLES AND ATTRIBUTES
% ============================================================

\section{Variables and Attributes}

\begin{frame}{Variables Data}
\small

Variables data are quantitative measurements.

Examples:

\begin{itemize}
    \item shaft diameter,
    \item thickness,
    \item mass,
    \item tensile strength,
    \item temperature,
    \item surface roughness.
\end{itemize}

A measurement may take many numerical values.

For example:

\[
24.98,\quad24.99,\quad25.01,\quad25.03\text{ mm}
\]

Variables charts generally use more information per observation.

\end{frame}

% ------------------------------------------------------------

\begin{frame}{Attributes Data}
\small

Attributes data classify observations.

Examples:

\begin{itemize}
    \item defective / non-defective,
    \item pass / fail,
    \item defect count,
    \item conforming / nonconforming.
\end{itemize}

Two important forms are:

\[
\boxed{
\text{Fraction defective}
}
\]

and

\[
\boxed{
\text{Number of defects}
}
\]

These lead to different control charts.

\end{frame}

% ------------------------------------------------------------

\begin{frame}{Variables vs Attributes}
\small

\begin{tabularx}{\textwidth}{YYY}
\toprule
Feature & Variables & Attributes\\
\midrule
Data type & Measured numerical value & Classification/count\\
Example & Diameter & Defective unit\\
Distribution & Often continuous & Binomial/Poisson models\\
Charts & $\bar X$, $R$, $S$ & $p$, $np$, $c$, $u$\\
Information & Detailed measurement & Classification/count\\
\bottomrule
\end{tabularx}

\end{frame}

% ============================================================
% SHEWHART CONTROL CHARTS
% ============================================================

\section{Shewhart Control Charts}

\begin{frame}{Basic Structure of a Control Chart}
\small

A Shewhart chart contains:

\[
\boxed{\text{Center Line}}
\]

\[
\boxed{\text{Upper Control Limit}}
\]

\[
\boxed{\text{Lower Control Limit}}
\]

For a generic statistic \(T\):

\[
UCL=\mu_T+3\sigma_T
\]

\[
CL=\mu_T
\]

\[
LCL=\mu_T-3\sigma_T
\]

approximately.

\end{frame}

% ------------------------------------------------------------

\begin{frame}{Why Three Sigma?}
\small

For a normal distribution:

\[
P(|Z|\le3)\approx0.9973.
\]

Therefore:

\[
P(|Z|>3)\approx0.0027.
\]

Thus approximately:

\[
99.73\%
\]

of observations are expected within \(\pm3\sigma\) when the
normal model is appropriate.

This provides a practical basis for Shewhart control limits.

\end{frame}

% ------------------------------------------------------------

\begin{frame}{Interpretation of a Control Chart}
\small

A process may show evidence of unusual behavior when observations
exhibit patterns such as:

\begin{itemize}
    \item a point beyond a control limit,
    \item a long run on one side of the center line,
    \item a sustained trend,
    \item systematic cycles,
    \item unusual clustering.
\end{itemize}

These patterns should be treated as signals for investigation.

\[
\boxed{
\text{Signal}\rightarrow\text{Investigate process}
}
\]

A control chart does not automatically identify the physical cause.

\end{frame}

% ============================================================
% XBAR-R CHART
% ============================================================

\section{$\bar X$--$R$ Charts}

\begin{frame}{$\bar X$--$R$ Control Charts}
\small

For measurable quality characteristics with rational subgroups,
the $\bar X$ and $R$ charts are commonly used together.

\[
\boxed{
\bar X\text{ chart}\rightarrow\text{process location}
}
\]

\[
\boxed{
R\text{ chart}\rightarrow\text{within-subgroup variation}
}
\]

A subgroup contains \(n\) observations.

For subgroup \(i\):

\[
\bar X_i=
\frac{1}{n}
\sum_{j=1}^{n}X_{ij}
\]

and:

\[
R_i=X_{\max}-X_{\min}.
\]

\end{frame}

% ------------------------------------------------------------

\begin{frame}{R Chart}
\small

For each subgroup:

\[
R_i=X_{\max,i}-X_{\min,i}.
\]

The average range is:

\[
\boxed{
\bar R=
\frac{1}{k}
\sum_{i=1}^{k}R_i
}
\]

where \(k\) is the number of subgroups.

Control limits:

\[
\boxed{
UCL_R=D_4\bar R
}
\]

\[
\boxed{
CL_R=\bar R
}
\]

\[
\boxed{
LCL_R=D_3\bar R
}
\]

The constants depend on subgroup size.

\end{frame}

% ------------------------------------------------------------

\begin{frame}{$\bar X$ Chart}
\small

Calculate the average of subgroup means:

\[
\boxed{
\bar{\bar X}
=
\frac{1}{k}
\sum_{i=1}^{k}\bar X_i
}
\]

Then:

\[
\boxed{
UCL_{\bar X}
=
\bar{\bar X}
+
A_2\bar R
}
\]

\[
\boxed{
CL_{\bar X}
=
\bar{\bar X}
}
\]

\[
\boxed{
LCL_{\bar X}
=
\bar{\bar X}
-
A_2\bar R
}
\]

The constant \(A_2\) depends on subgroup size.

\end{frame}

% ------------------------------------------------------------

\begin{frame}{Why Rational Subgroups?}
\small

A rational subgroup should contain observations collected under
conditions where within-subgroup variation represents short-term
process variation.

Between-subgroup differences can then reveal:

\begin{itemize}
    \item shifts in process mean,
    \item changes in process conditions,
    \item gradual trends,
    \item special causes.
\end{itemize}

Poor subgroup construction can hide important process behavior.

\end{frame}

% ------------------------------------------------------------

\begin{frame}{Important $\bar X$--$R$ Constants}
\small

Typical constants include:

\begin{tabularx}{\textwidth}{CCCCC}
\toprule
\(n\) & \(A_2\) & \(D_3\) & \(D_4\) & \(d_2\)\\
\midrule
2 & 1.880 & 0 & 3.267 & 1.128\\
3 & 1.023 & 0 & 2.574 & 1.693\\
4 & 0.729 & 0 & 2.282 & 2.059\\
5 & 0.577 & 0 & 2.114 & 2.326\\
6 & 0.483 & 0 & 2.004 & 2.534\\
\bottomrule
\end{tabularx}

The constants depend on the statistical relationship between
range and process standard deviation.

\end{frame}

% ------------------------------------------------------------

\begin{frame}{Estimating Process Standard Deviation}
\small

For a range-based estimate:

\[
\boxed{
\hat\sigma=
\frac{\bar R}{d_2}
}
\]

where \(d_2\) depends on subgroup size.

This relationship is based on the expected range of a sample
from a normal population.

Therefore the range can provide an estimate of short-term
process dispersion.

\end{frame}

% ------------------------------------------------------------

\begin{frame}{$\bar X$--$R$: Numerical Example}
\small

Suppose five subgroups of size four have:

\[
\bar X_i=
10.2,\ 10.4,\ 10.1,\ 10.3,\ 10.5
\]

and:

\[
R_i=
0.4,\ 0.5,\ 0.3,\ 0.4,\ 0.6.
\]

Then:

\[
\bar{\bar X}
=
\frac{10.2+10.4+10.1+10.3+10.5}{5}
\]

\[
\boxed{
\bar{\bar X}=10.30
}
\]

and:

\[
\boxed{
\bar R=\frac{0.4+0.5+0.3+0.4+0.6}{5}=0.44
}
\]

\end{frame}

% ------------------------------------------------------------

\begin{frame}{$\bar X$--$R$: Control Limits}
\small

For \(n=4\):

\[
A_2=0.729
\]

\[
D_3=0
\]

\[
D_4=2.282.
\]

Therefore:

\[
UCL_{\bar X}
=
10.30+(0.729)(0.44)
\]

\[
\boxed{
UCL_{\bar X}=10.621
}
\]

and:

\[
LCL_{\bar X}
=
10.30-(0.729)(0.44)
\]

\[
\boxed{
LCL_{\bar X}=9.979
}
\]

\end{frame}

% ------------------------------------------------------------

\begin{frame}{R Chart: Numerical Calculation}
\small

Using:

\[
\bar R=0.44
\]

and:

\[
D_3=0,\qquad D_4=2.282,
\]

we obtain:

\[
UCL_R=(2.282)(0.44)
\]

\[
\boxed{
UCL_R=1.004
}
\]

and:

\[
LCL_R=(0)(0.44)=0.
\]

Thus:

\[
\boxed{
0\le R\le1.004
}
\]

represents the approximate control region for the range chart.

\end{frame}

% ============================================================
% DEFECT / DEFECTIVE
% ============================================================

\section{Defects and Defectives}

\begin{frame}{Defect vs Defective}
\small

A \textbf{defect} is a specific nonconformity.

A \textbf{defective unit} is a unit having one or more defects that
cause it to be classified as nonconforming.

Example:

One component may contain:

\[
\text{Scratch}+\text{Dent}+\text{Dimensional error}
\]

which represents:

\[
3\text{ defects}
\]

but potentially only:

\[
1\text{ defective unit}.
\]

\end{frame}

% ============================================================
% P CHART
% ============================================================

\section{$p$ Chart}

\begin{frame}{$p$ Control Chart}
\small

The \(p\)-chart monitors the fraction or proportion of defective
units in a sample.

For sample \(i\):

\[
\boxed{
p_i=
\frac{d_i}{n_i}
}
\]

where:

\begin{itemize}
    \item \(d_i\) = number of defective units,
    \item \(n_i\) = sample size.
\end{itemize}

If the process fraction defective is \(p\):

\[
D_i\sim\mathrm{Binomial}(n_i,p).
\]

\end{frame}

% ------------------------------------------------------------

\begin{frame}{Mean Fraction Defective}
\small

For equal or varying sample sizes, an estimate of the overall
fraction defective is:

\[
\boxed{
\bar p
=
\frac{\sum d_i}{\sum n_i}
}
\]

Then:

\[
\boxed{
\bar q=1-\bar p
}
\]

For equal sample sizes this is equivalent to averaging the
sample proportions.

\end{frame}

% ------------------------------------------------------------

\begin{frame}{$p$ Chart Control Limits}
\small

For sample size \(n\):

\[
\sigma_p=
\sqrt{
\frac{\bar p(1-\bar p)}{n}
}
\]

Therefore:

\[
\boxed{
UCL_p=
\bar p+
3
\sqrt{
\frac{\bar p(1-\bar p)}{n}
}
}
\]

\[
\boxed{
CL_p=\bar p
}
\]

\[
\boxed{
LCL_p=
\bar p-
3
\sqrt{
\frac{\bar p(1-\bar p)}{n}
}
}
\]

If \(LCL_p<0\), it is set to zero.

\end{frame}

% ------------------------------------------------------------

\begin{frame}{$p$ Chart: Numerical Example}
\small

Suppose:

\[
n=100
\]

units are inspected in each sample.

Observed defectives:

\[
8,\ 5,\ 7,\ 10,\ 6.
\]

Total defectives:

\[
\sum d_i=36.
\]

Total inspected:

\[
5(100)=500.
\]

Therefore:

\[
\bar p=\frac{36}{500}=0.072.
\]

\[
\boxed{
\bar p=7.2\%
}
\]

\end{frame}

% ------------------------------------------------------------

\begin{frame}{$p$ Chart: Control Limits}
\small

For:

\[
\bar p=0.072,\qquad n=100,
\]

the standard deviation is:

\[
\sigma_p
=
\sqrt{
\frac{0.072(0.928)}{100}
}
\]

\[
\sigma_p\approx0.02585.
\]

Therefore:

\[
UCL=0.072+3(0.02585)
\]

\[
\boxed{
UCL\approx0.1496
}
\]

and:

\[
LCL=0.072-3(0.02585)<0.
\]

Hence:

\[
\boxed{
LCL=0
}
\]

\end{frame}

% ------------------------------------------------------------

\begin{frame}{Unequal Sample Sizes in a $p$ Chart}
\small

If sample size varies:

\[
n_1\neq n_2\neq\cdots
\]

then each sample has a different standard error.

For sample \(i\):

\[
\boxed{
UCL_i=
\bar p+
3
\sqrt{
\frac{\bar p(1-\bar p)}{n_i}
}
}
\]

\[
\boxed{
LCL_i=
\bar p-
3
\sqrt{
\frac{\bar p(1-\bar p)}{n_i}
}
}
\]

Thus the control limits vary with sample size.

\end{frame}

% ============================================================
% C CHART
% ============================================================

\section{$c$ Chart}

\begin{frame}{$c$ Control Chart}
\small

The \(c\)-chart monitors the number of defects per inspection unit
when the opportunity for defects is constant.

The usual model is:

\[
\boxed{
C\sim\mathrm{Poisson}(c)
}
\]

For a Poisson distribution:

\[
E(C)=c
\]

and:

\[
Var(C)=c.
\]

Therefore:

\[
\sigma_c=\sqrt c.
\]

\end{frame}

% ------------------------------------------------------------

\begin{frame}{$c$ Chart Control Limits}
\small

If the average number of defects is:

\[
\bar c=
\frac{1}{k}\sum_{i=1}^{k}c_i,
\]

then:

\[
\boxed{
UCL_c=\bar c+3\sqrt{\bar c}
}
\]

\[
\boxed{
CL_c=\bar c
}
\]

\[
\boxed{
LCL_c=\bar c-3\sqrt{\bar c}
}
\]

If the lower limit is negative:

\[
\boxed{
LCL_c=0
}
\]

\end{frame}

% ------------------------------------------------------------

\begin{frame}{$c$ Chart: Numerical Example}
\small

Suppose the observed defect counts are:

\[
4,\ 7,\ 5,\ 6,\ 8.
\]

Then:

\[
\bar c=
\frac{4+7+5+6+8}{5}
\]

\[
\boxed{
\bar c=6
}
\]

The standard deviation is:

\[
\sqrt 6=2.449.
\]

Therefore:

\[
UCL=6+3(2.449)
\]

\[
\boxed{
UCL\approx13.35
}
\]

\end{frame}

% ------------------------------------------------------------

\begin{frame}{$p$ Chart vs $c$ Chart}
\small

\begin{tabularx}{\textwidth}{YYY}
\toprule
Feature & $p$ chart & $c$ chart\\
\midrule
Quantity & Fraction defective & Number of defects\\
Unit status & Defective/nondefective & Multiple defects possible\\
Distribution & Binomial & Poisson\\
Sample size & Can vary & Inspection opportunity normally constant\\
Example & Defective shafts & Scratches per sheet\\
\bottomrule
\end{tabularx}

\vspace{6pt}

The choice depends on what is being counted.

\end{frame}

% ============================================================
% ACCEPTANCE SAMPLING
% ============================================================

\section{Acceptance Sampling}

\begin{frame}{What is Acceptance Sampling?}
\small

Acceptance sampling is a procedure in which a sample is taken
from a lot and the lot is accepted or rejected based on sample
results.

It is different from process control.

\[
\boxed{
\text{Control chart}
\rightarrow
\text{controls a process}
}
\]

\[
\boxed{
\text{Acceptance sampling}
\rightarrow
\text{makes a lot decision}
}
\]

\end{frame}

% ------------------------------------------------------------

\begin{frame}{Why Use Acceptance Sampling?}
\small

100\% inspection may be undesirable when:

\begin{itemize}
    \item inspection is destructive,
    \item inspection is expensive,
    \item the lot is very large,
    \item inspection time is excessive,
    \item only a statistical decision is required.
\end{itemize}

However, sampling introduces risks because the entire lot
is not observed.

\end{frame}

% ------------------------------------------------------------

\begin{frame}{Single Sampling Plan}
\small

A single sampling plan is defined by:

\[
\boxed{
(n,c)
}
\]

where:

\begin{itemize}
    \item \(n\) = sample size,
    \item \(c\) = acceptance number.
\end{itemize}

Procedure:

\begin{enumerate}
    \item Select \(n\) units.
    \item Count nonconforming units.
    \item If number defective \(\le c\), accept the lot.
    \item If number defective \(>c\), reject the lot.
\end{enumerate}

\end{frame}

% ------------------------------------------------------------

\begin{frame}{Binomial Model for Acceptance Sampling}
\small

For a large lot or independent sampling approximation:

\[
X\sim\mathrm{Binomial}(n,p)
\]

where:

\[
P(X=x)
=
{n\choose x}
p^x(1-p)^{n-x}.
\]

For a single sampling plan \((n,c)\):

\[
\boxed{
P_a(p)
=
P(X\le c)
}
\]

Therefore:

\[
\boxed{
P_a(p)
=
\sum_{x=0}^{c}
{n\choose x}
p^x(1-p)^{n-x}
}
\]

\end{frame}

% ------------------------------------------------------------

\begin{frame}{Acceptance Probability Example}
\small

Consider:

\[
n=20,\qquad c=1.
\]

If:

\[
p=0.05,
\]

the lot is accepted when:

\[
X=0
\]

or:

\[
X=1.
\]

Thus:

\[
P_a=P(X=0)+P(X=1).
\]

For the binomial distribution:

\[
P(X=0)=(0.95)^{20}
\]

and:

\[
P(X=1)=20(0.05)(0.95)^{19}.
\]

\end{frame}

% ------------------------------------------------------------

\begin{frame}{Numerical Acceptance Probability}
\small

For:

\[
n=20,\quad c=1,\quad p=0.05
\]

we obtain approximately:

\[
P(X=0)\approx0.3585
\]

and:

\[
P(X=1)\approx0.3774.
\]

Therefore:

\[
P_a
\approx
0.3585+0.3774
\]

\[
\boxed{
P_a\approx0.7359
}
\]

Thus approximately \(73.6\%\) of such lots would be accepted
under this model.

\end{frame}

% ------------------------------------------------------------

\begin{frame}{Hypergeometric Model}
\small

For sampling without replacement from a finite lot, the
hypergeometric distribution may be appropriate.

If:

\begin{itemize}
    \item \(N\) = lot size,
    \item \(D\) = defective units in lot,
    \item \(n\) = sample size,
    \item \(x\) = defectives in sample,
\end{itemize}

then:

\[
\boxed{
P(X=x)
=
\frac{
{D\choose x}
{N-D\choose n-x}
}{
{N\choose n}
}
}
\]

This accounts explicitly for finite-population sampling.

\end{frame}

% ============================================================
% DOUBLE SAMPLING
% ============================================================

\section{Double Sampling}

\begin{frame}{Double Sampling Plan}
\small

A double sampling plan allows a second sample when the first
sample does not provide a sufficiently clear decision.

Typical sequence:

\[
\boxed{
\text{First sample}
\rightarrow
\begin{cases}
\text{Accept}\\
\text{Reject}\\
\text{Take second sample}
\end{cases}
}
\]

The second sample can reduce inspection effort for lots that
are clearly good or clearly bad.

\end{frame}

% ------------------------------------------------------------

\begin{frame}{Conceptual Double Sampling Example}
\small

Suppose the first sample contains:

\[
n_1=50
\]

units.

Possible decision:

\begin{itemize}
    \item If defectives \(\le c_1\): accept.
    \item If defectives \(>r_1\): reject.
    \item If \(c_1<d_1\le r_1\): take second sample.
\end{itemize}

After the second sample:

\[
d_1+d_2
\]

is compared with a second acceptance criterion.

\vspace{5pt}

The exact plan is determined by the specified sampling standard.

\end{frame}

% ============================================================
% OC CURVES
% ============================================================

\section{Operating Characteristic Curves}

\begin{frame}{Operating Characteristic Curve}
\small

The \textbf{Operating Characteristic (OC) curve} shows:

\[
\boxed{
P_a
\text{ versus }
p
}
\]

where:

\begin{itemize}
    \item \(P_a\) = probability of accepting the lot,
    \item \(p\) = fraction defective in the lot.
\end{itemize}

For a fixed sampling plan:

\[
(n,c),
\]

the OC curve describes how the plan behaves for different
quality levels.

\end{frame}

% ------------------------------------------------------------

\begin{frame}{Shape of an OC Curve}
\small

A typical OC curve has:

\[
P_a\approx1
\]

for very low defect levels.

As the lot defect level increases:

\[
P_a\downarrow
\]

because it becomes increasingly difficult for the lot to satisfy
the acceptance criterion.

Therefore:

\[
\boxed{
\text{Higher lot defect level}
\rightarrow
\text{lower acceptance probability}
}
\]

for a conventional single sampling plan.

\end{frame}

% ------------------------------------------------------------

\begin{frame}{Effect of Acceptance Number \(c\)}
\small

For fixed \(n\):

\begin{itemize}
    \item Increasing \(c\) makes acceptance possible for more
    defective units.
    \item Decreasing \(c\) makes the acceptance criterion stricter.
\end{itemize}

The acceptance probability is:

\[
P_a=P(X\le c).
\]

Therefore changing \(c\) directly changes the OC curve.

\end{frame}

% ------------------------------------------------------------

\begin{frame}{Effect of Sample Size \(n\)}
\small

For a fixed acceptance rule, increasing sample size generally
provides more information about the lot.

A larger sample can distinguish different quality levels more
sharply.

However:

\[
\boxed{
\text{larger }n
\rightarrow
\text{higher inspection cost}
}
\]

Therefore sampling-plan design involves balancing information
and inspection effort.

\end{frame}

% ============================================================
% AQL / LTPD / RISKS
% ============================================================

\section{AQL, LTPD and Risks}

\begin{frame}{Acceptable Quality Level}
\small

The \textbf{AQL} is a specified quality level used in sampling-plan
design and quality assurance.

It is associated with a high probability of acceptance under the
chosen sampling plan.

It should not be interpreted as:

\[
\boxed{
\text{AQL = guaranteed defect level}
}
\]

Instead, it is a parameter within a statistical sampling system.

\end{frame}

% ------------------------------------------------------------

\begin{frame}{LTPD / RQL}
\small

The \textbf{Lot Tolerance Percent Defective (LTPD)} or related
Rejectable Quality Level (RQL) represents an undesirable quality
level used to characterize consumer protection.

At this quality level, the probability of accepting the lot is
designed to be relatively low.

Thus:

\[
\boxed{
\text{AQL}\rightarrow\text{producer-side quality point}
}
\]

\[
\boxed{
\text{LTPD/RQL}\rightarrow\text{consumer-side protection point}
}
\]

Exact definitions depend on the sampling standard being used.

\end{frame}

% ------------------------------------------------------------

\begin{frame}{Producer's Risk \(\alpha\)}
\small

Producer's risk is the probability of rejecting a lot considered
to be of acceptable quality.

At the chosen AQL:

\[
\boxed{
\alpha
=
P(\text{reject}\mid\text{acceptable quality})
}
\]

Since:

\[
P_a=P(\text{accept}),
\]

we have:

\[
\boxed{
\alpha=1-P_a
}
\]

at the specified acceptable-quality point.

\end{frame}

% ------------------------------------------------------------

\begin{frame}{Consumer's Risk \(\beta\)}
\small

Consumer's risk is the probability of accepting a lot associated
with an undesirable quality level.

At the specified LTPD/RQL:

\[
\boxed{
\beta
=
P(\text{accept}\mid\text{poor quality})
}
\]

Therefore:

\[
\boxed{
\beta=P_a(p_{\mathrm{LTPD}})
}
\]

The sampling plan is designed to control both types of risk
according to the specified requirements.

\end{frame}

% ------------------------------------------------------------

\begin{frame}{Producer vs Consumer Risk}
\small

\begin{tabularx}{\textwidth}{YYY}
\toprule
Feature & Producer's risk & Consumer's risk\\
\midrule
Symbol & $\alpha$ & $\beta$\\
Concern & Good lot rejected & Poor lot accepted\\
Quality point & AQL & LTPD/RQL\\
Decision error & False rejection & False acceptance\\
\bottomrule
\end{tabularx}

\vspace{6pt}

Both risks exist because only a sample is inspected.

\end{frame}

% ============================================================
% CONTROL CHART VS ACCEPTANCE SAMPLING
% ============================================================

\section{Control Charts vs Acceptance Sampling}

\begin{frame}{Process Control vs Lot Acceptance}
\small

\begin{tabularx}{\textwidth}{YYY}
\toprule
Feature & Control chart & Acceptance sampling\\
\midrule
Purpose & Monitor process & Decide about lot\\
Data sequence & Sequential & Sample-based\\
Main question & Is process stable? & Accept or reject lot?\\
Typical tool & $\bar X$, $R$, $p$, $c$ & Single/double sampling\\
Action & Investigate process & Lot decision\\
Primary focus & Process variation & Lot quality\\
\bottomrule
\end{tabularx}

\end{frame}

% ============================================================
% ASSUMPTIONS
% ============================================================

\section{Assumptions and Limitations}

\begin{frame}{Assumptions in Time Study}
\small

Typical assumptions include:

\begin{itemize}
    \item method is clearly defined,
    \item worker is qualified,
    \item observations are representative,
    \item equipment is operating normally,
    \item abnormal interruptions are identified,
    \item rating is applied consistently,
    \item allowance definitions are clear.
\end{itemize}

\vspace{5pt}

If these assumptions are violated, the resulting standard time may
not represent the intended operating condition.

\end{frame}

% ------------------------------------------------------------

\begin{frame}{Limitations of Time Study}
\small

Potential limitations:

\begin{itemize}
    \item observer influence,
    \item rating subjectivity,
    \item insufficient observations,
    \item worker-to-worker differences,
    \item method changes during study,
    \item unusual operating conditions,
    \item stopwatch resolution,
    \item fatigue effects.
\end{itemize}

Therefore a standard time should be viewed as a defined engineering
standard, not as an immutable physical constant.

\end{frame}

% ------------------------------------------------------------

\begin{frame}{Assumptions in Work Sampling}
\small

Work sampling relies on:

\begin{itemize}
    \item representative observations,
    \item appropriate randomization,
    \item sufficiently large sample size,
    \item independent or approximately independent observations,
    \item clearly defined activity categories.
\end{itemize}

Poor randomization can create systematic sampling bias.

For example, observing only during morning hours may not represent
the entire workday.

\end{frame}

% ------------------------------------------------------------

\begin{frame}{Assumptions in Control Charts}
\small

Depending on the chart, assumptions may include:

\begin{itemize}
    \item observations are appropriately subgrouped,
    \item measurement system is adequate,
    \item process conditions are reasonably comparable,
    \item the selected probability model is suitable,
    \item sampling frequency is appropriate.
\end{itemize}

Control limits should be recalculated when the underlying process
or sampling structure changes materially.

\end{frame}

% ============================================================
% INTEGRATED MANUFACTURING EXAMPLE
% ============================================================

\section{Integrated Engineering Example}

\begin{frame}{Integrated Manufacturing Scenario}
\small

Consider a shaft-manufacturing line.

The process consists of:

\[
\text{Raw material}
\rightarrow
\text{Turning}
\rightarrow
\text{Grinding}
\rightarrow
\text{Inspection}
\rightarrow
\text{Packing}
\]

Possible engineering questions:

\begin{itemize}
    \item Is the method efficient?
    \item What is the standard machining time?
    \item How much operator idle time occurs?
    \item Is shaft diameter statistically stable?
    \item What fraction of shafts are defective?
    \item Should incoming lots be accepted?
\end{itemize}

\end{frame}

% ------------------------------------------------------------

\begin{frame}{Integrated Method Study}
\small

The method-study sequence may identify:

\begin{itemize}
    \item unnecessary material movement,
    \item repeated inspection,
    \item poor workstation arrangement,
    \item excessive operator travel,
    \item unnecessary handling.
\end{itemize}

ECRS can then be applied:

\[
\boxed{
\text{Eliminate}
\rightarrow
\text{Combine}
\rightarrow
\text{Rearrange}
\rightarrow
\text{Simplify}
}
\]

The improved method should be defined before establishing
the final standard time.

\end{frame}

% ------------------------------------------------------------

\begin{frame}{Integrated Time Study}
\small

Suppose:

\[
OT=3.20\text{ min}
\]

and:

\[
R=105\%.
\]

Then:

\[
NT=3.20(1.05)=3.36\text{ min}.
\]

If allowance is 15\% of normal time:

\[
ST=3.36(1.15)
\]

\[
\boxed{
ST=3.864\text{ min}
}
\]

Thus the production standard is approximately:

\[
\boxed{
3.86\text{ min/unit}
}
\]

\end{frame}

% ------------------------------------------------------------

\begin{frame}{Capacity Interpretation}
\small

If:

\[
ST=3.864\text{ min/unit},
\]

then theoretical output per 60-minute hour is:

\[
Q=
\frac{60}{3.864}
\]

\[
\boxed{
Q\approx15.52\text{ units/hour}
}
\]

This does not automatically mean 15.52 saleable units will be
produced because actual output depends on:

\begin{itemize}
    \item downtime,
    \item quality losses,
    \item material availability,
    \item setup,
    \item maintenance,
    \item scheduling.
\end{itemize}

\end{frame}

% ------------------------------------------------------------

\begin{frame}{Integrated Quality Monitoring}
\small

Suppose shaft diameter is measured in rational subgroups.

Use:

\[
\boxed{\bar X\text{ chart}}
\]

to monitor process location and:

\[
\boxed{R\text{ chart}}
\]

to monitor within-subgroup variation.

If the product is classified as defective/nondefective:

\[
\boxed{p\text{ chart}}
\]

can monitor the fraction defective.

Thus several statistical tools may be used simultaneously.

\end{frame}

% ============================================================
% FORMULA SHEET
% ============================================================

\section{Formula Sheet}

\begin{frame}{Formula Sheet: Time Study}
\small

\[
\boxed{
OT=\frac{\sum T_i}{n}
}
\]

\[
\boxed{
NT=OT\left(\frac{R}{100}\right)
}
\]

\[
\boxed{
ST=NT(1+A)
}
\]

when allowance is a fraction of normal time.

\[
\boxed{
ST=\frac{NT}{1-A}
}
\]

when allowance is a fraction of standard time.

\end{frame}

% ------------------------------------------------------------

\begin{frame}{Formula Sheet: Work Sampling}
\small

\[
\boxed{
\hat p=\frac{x}{n}
}
\]

\[
\boxed{
SE(\hat p)
=
\sqrt{
\frac{\hat p(1-\hat p)}{n}
}
}
\]

\[
\boxed{
CI=
\hat p
\pm
z_{\alpha/2}
SE(\hat p)
}
\]

\[
\boxed{
n=
\frac{
z_{\alpha/2}^{2}pq
}{
e^2
}
}
\]

For maximum variance:

\[
\boxed{
p=q=0.5
}
\]

\end{frame}

% ------------------------------------------------------------

\begin{frame}{Formula Sheet: $\bar X$--$R$}
\small

\[
\boxed{
\bar X_i=
\frac{1}{n}
\sum X_{ij}
}
\]

\[
\boxed{
R_i=X_{\max}-X_{\min}
}
\]

\[
\boxed{
\bar{\bar X}
=
\frac{1}{k}
\sum\bar X_i
}
\]

\[
\boxed{
\bar R=
\frac{1}{k}
\sum R_i
}
\]

\[
\boxed{
UCL_{\bar X}
=
\bar{\bar X}+A_2\bar R
}
\]

\[
\boxed{
LCL_{\bar X}
=
\bar{\bar X}-A_2\bar R
}
\]

\end{frame}

% ------------------------------------------------------------

\begin{frame}{Formula Sheet: R Chart}
\small

\[
\boxed{
UCL_R=D_4\bar R
}
\]

\[
\boxed{
CL_R=\bar R
}
\]

\[
\boxed{
LCL_R=D_3\bar R
}
\]

Process standard deviation estimate:

\[
\boxed{
\hat\sigma=\frac{\bar R}{d_2}
}
\]

The constants:

\[
A_2,\quad D_3,\quad D_4,\quad d_2
\]

depend on subgroup size.

\end{frame}

% ------------------------------------------------------------

\begin{frame}{Formula Sheet: $p$ Chart}
\small

\[
\boxed{
p_i=\frac{d_i}{n_i}
}
\]

\[
\boxed{
\bar p=
\frac{\sum d_i}{\sum n_i}
}
\]

For constant \(n\):

\[
\boxed{
UCL_p=
\bar p+
3\sqrt{
\frac{\bar p(1-\bar p)}{n}
}
}
\]

\[
\boxed{
LCL_p=
\bar p-
3\sqrt{
\frac{\bar p(1-\bar p)}{n}
}
}
\]

with negative \(LCL\) replaced by zero.

\end{frame}

% ------------------------------------------------------------

\begin{frame}{Formula Sheet: $c$ Chart}
\small

\[
\boxed{
\bar c=
\frac{1}{k}\sum c_i
}
\]

\[
\boxed{
UCL_c=
\bar c+3\sqrt{\bar c}
}
\]

\[
\boxed{
CL_c=\bar c
}
\]

\[
\boxed{
LCL_c=
\bar c-3\sqrt{\bar c}
}
\]

If:

\[
LCL_c<0,
\]

use:

\[
\boxed{
LCL_c=0.
}
\]

\end{frame}

% ------------------------------------------------------------

\begin{frame}{Formula Sheet: Acceptance Sampling}
\small

Binomial probability:

\[
\boxed{
P(X=x)
=
{n\choose x}
p^x(1-p)^{n-x}
}
\]

Acceptance probability:

\[
\boxed{
P_a(p)=
\sum_{x=0}^{c}
{n\choose x}
p^x(1-p)^{n-x}
}
\]

Hypergeometric probability:

\[
\boxed{
P(X=x)
=
\frac{
{D\choose x}
{N-D\choose n-x}
}{
{N\choose n}
}
}
\]

\end{frame}

% ============================================================
% COMPARISON TABLES
% ============================================================

\section{Comparison Tables}

\begin{frame}{Method Study vs Work Measurement}
\small

\begin{tabularx}{\textwidth}{YYY}
\toprule
Feature & Method Study & Work Measurement\\
\midrule
Main question & How should work be done? & How long should work take?\\
Primary focus & Method & Time\\
Main output & Improved method & Standard time\\
Typical tools & Charts, ECRS & Stopwatch, sampling\\
Sequence & Usually before measurement & After method definition\\
\bottomrule
\end{tabularx}

\end{frame}

% ------------------------------------------------------------

\begin{frame}{Time Study vs Work Sampling}
\small

\begin{tabularx}{\textwidth}{YYY}
\toprule
Feature & Time Study & Work Sampling\\
\midrule
Observation & Continuous/repetitive & Random observations\\
Best suited for & Short repetitive cycles & Irregular activities\\
Output & Task time & Proportion of time\\
Equipment & Stopwatch/timer & Sampling schedule\\
Statistical basis & Direct measurement & Probability sampling\\
\bottomrule
\end{tabularx}

\end{frame}

% ------------------------------------------------------------

\begin{frame}{Control Charts Comparison}
\small

\begin{tabularx}{\textwidth}{YYYY}
\toprule
Chart & Data & Model & Main purpose\\
\midrule
$\bar X$ & Mean & Normal approximation & Location\\
$R$ & Range & Range distribution & Dispersion\\
$p$ & Fraction defective & Binomial & Defective fraction\\
$c$ & Defect count & Poisson & Defect count\\
\bottomrule
\end{tabularx}

\end{frame}

% ------------------------------------------------------------

\begin{frame}{Important Distinctions}
\small

\[
\boxed{
\text{Defect}\neq\text{Defective}
}
\]

\[
\boxed{
\text{Control limit}\neq\text{Specification limit}
}
\]

\[
\boxed{
\text{Normal time}\neq\text{Standard time}
}
\]

\[
\boxed{
\text{Process control}\neq\text{Acceptance sampling}
}
\]

\[
\boxed{
\text{Variables data}\neq\text{Attributes data}
}
\]

These distinctions are frequently tested in examinations and viva.

\end{frame}

% ============================================================
% DETAILED CONCEPTUAL QUESTIONS
% ============================================================

\section{Viva and Examination Questions}

\begin{frame}{Viva Questions: Work Study}
\small

\begin{enumerate}
    \item What is work study?
    \item What are the two main components of work study?
    \item Why is method study performed before time study?
    \item What is the SREDIM procedure?
    \item What is ECRS?
    \item What is the purpose of an operation process chart?
    \item What is a flow process chart?
    \item Why is a two-handed chart useful?
    \item What information is provided by a man-machine chart?
\end{enumerate}

\end{frame}

% ------------------------------------------------------------

\begin{frame}{Viva Questions: Time Study}
\small

\begin{enumerate}
    \item What is work measurement?
    \item Why is a job divided into elements?
    \item What is observed time?
    \item What is normal time?
    \item Why is performance rating required?
    \item What does a rating of 100\% mean?
    \item What happens when rating exceeds 100\%?
    \item Why are allowances added?
    \item What is standard time?
    \item Why can two allowance formulas exist?
\end{enumerate}

\end{frame}

% ------------------------------------------------------------

\begin{frame}{Viva Questions: Work Sampling}
\small

\begin{enumerate}
    \item What is work sampling?
    \item What is the probability model behind work sampling?
    \item Why is random sampling important?
    \item What is \(\hat p\)?
    \item What is the standard error of a proportion?
    \item How does sample size affect precision?
    \item Why is \(p=0.5\) conservative for sample-size planning?
    \item What is a confidence interval?
\end{enumerate}

\end{frame}

% ------------------------------------------------------------

\begin{frame}{Viva Questions: SQC}
\small

\begin{enumerate}
    \item What is statistical quality control?
    \item What is process variation?
    \item What are common causes?
    \item What are special causes?
    \item What is a control chart?
    \item Why are three-sigma limits commonly used?
    \item What is the difference between control and specification limits?
    \item What is rational subgrouping?
\end{enumerate}

\end{frame}

% ------------------------------------------------------------

\begin{frame}{Viva Questions: Control Charts}
\small

\begin{enumerate}
    \item What is an $\bar X$ chart?
    \item What is an \(R\) chart?
    \item Why are $\bar X$ and \(R\) charts used together?
    \item What is \(A_2\)?
    \item What is \(D_3\)?
    \item What is \(D_4\)?
    \item What is \(d_2\)?
    \item When is a \(p\)-chart used?
    \item When is a \(c\)-chart used?
\end{enumerate}

\end{frame}

% ------------------------------------------------------------

\begin{frame}{Viva Questions: Acceptance Sampling}
\small

\begin{enumerate}
    \item What is acceptance sampling?
    \item What is a single sampling plan?
    \item What do \(n\) and \(c\) represent?
    \item What is acceptance probability?
    \item When is the binomial model appropriate?
    \item When is the hypergeometric model appropriate?
    \item What is an OC curve?
    \item What is AQL?
    \item What is LTPD?
    \item What are producer's and consumer's risks?
\end{enumerate}

\end{frame}

% ============================================================
% NUMERICAL PRACTICE
% ============================================================

\section{Numerical Practice}

\begin{frame}{Numerical Problem 1: Normal and Standard Time}
\small

A job has an observed time of:

\[
OT=4.5\text{ min}.
\]

The worker rating is:

\[
R=110\%.
\]

The allowance is:

\[
A=20\%
\]

of normal time.

Find:

\begin{enumerate}
    \item Normal time
    \item Standard time
\end{enumerate}

\vspace{5pt}

\textbf{Solution:}

\[
NT=4.5(1.10)=4.95
\]

\[
ST=4.95(1.20)=5.94
\]

\[
\boxed{
ST=5.94\text{ min}
}
\]

\end{frame}

% ------------------------------------------------------------

\begin{frame}{Numerical Problem 2: Work Sampling}
\small

During 800 random observations, a machine was found to be idle
in 120 observations.

Find:

\begin{enumerate}
    \item Estimated idle proportion
    \item Standard error
    \item Approximate 95\% confidence interval
\end{enumerate}

\vspace{4pt}

\[
\hat p=\frac{120}{800}=0.15
\]

\[
SE=
\sqrt{
\frac{0.15(0.85)}{800}
}
\approx0.01263
\]

Therefore:

\[
CI=
0.15\pm1.96(0.01263)
\]

\[
\boxed{
CI\approx(0.1252,\ 0.1748)
}
\]

\end{frame}

% ------------------------------------------------------------

\begin{frame}{Numerical Problem 3: $\bar X$ Chart}
\small

Given:

\[
\bar{\bar X}=50.20
\]

\[
\bar R=0.80
\]

and for subgroup size \(n=5\):

\[
A_2=0.577.
\]

Find control limits.

\vspace{4pt}

\[
UCL=
50.20+(0.577)(0.80)
\]

\[
\boxed{
UCL=50.6616
}
\]

\[
LCL=
50.20-(0.577)(0.80)
\]

\[
\boxed{
LCL=49.7384
}
\]

\end{frame}

% ------------------------------------------------------------

\begin{frame}{Numerical Problem 4: $p$ Chart}
\small

Suppose:

\[
\bar p=0.04
\]

and:

\[
n=200.
\]

Then:

\[
\sigma_p=
\sqrt{
\frac{0.04(0.96)}{200}
}
\]

\[
\sigma_p\approx0.01386.
\]

Hence:

\[
UCL=0.04+3(0.01386)
\]

\[
\boxed{
UCL\approx0.0816
}
\]

and:

\[
LCL=0.04-3(0.01386)<0.
\]

Therefore:

\[
\boxed{
LCL=0
}
\]

\end{frame}

% ------------------------------------------------------------

\begin{frame}{Numerical Problem 5: $c$ Chart}
\small

Suppose:

\[
\bar c=9.
\]

Then:

\[
\sigma_c=\sqrt9=3.
\]

Therefore:

\[
UCL=9+3(3)
\]

\[
\boxed{
UCL=18
}
\]

and:

\[
LCL=9-3(3)=0.
\]

Hence:

\[
\boxed{
0\le c\le18
}
\]

is the approximate three-sigma control region.

\end{frame}

% ============================================================
% DEEP CONCEPTUAL INTERPRETATION
% ============================================================

\section{Physical and Engineering Interpretation}

\begin{frame}{Physical Meaning of Standard Time}
\small

Standard time is not simply the stopwatch reading.

Conceptually:

\[
\boxed{
ST=
\text{time at standard performance}
+
\text{appropriate allowances}
}
\]

It represents the expected time under defined conditions.

Therefore standard time can support:

\begin{itemize}
    \item production scheduling,
    \item labour planning,
    \item machine loading,
    \item costing,
    \item productivity analysis.
\end{itemize}

\end{frame}

% ------------------------------------------------------------

\begin{frame}{Physical Meaning of Control Limits}
\small

Control limits describe the expected statistical behavior of a
process under the assumed model.

They are not engineering tolerances.

For example:

\[
\text{Diameter target}=25.00\text{ mm}
\]

might have specifications:

\[
24.95\le D\le25.05\text{ mm}.
\]

Control limits are instead calculated from observed process
variation.

\[
\boxed{
\text{Specification = requirement}
}
\]

\[
\boxed{
\text{Control limit = process behavior}
}
\]

\end{frame}

% ------------------------------------------------------------

\begin{frame}{Physical Meaning of Process Capability}
\small

A stable process can still produce excessive nonconforming units
if its natural variation is large relative to the specifications.

Conceptually:

\[
\boxed{
\text{Stability}
\neq
\text{Capability}
}
\]

Stability asks:

\[
\text{Is the process statistically predictable?}
\]

Capability asks:

\[
\text{Can the predictable process meet specifications?}
\]

Both questions are important.

\end{frame}

% ------------------------------------------------------------

\begin{frame}{Why Statistical Thinking Matters}
\small

Engineering decisions should distinguish:

\[
\text{signal}
\]

from:

\[
\text{noise}.
\]

A single unusual observation may be:

\begin{itemize}
    \item a genuine special cause,
    \item measurement error,
    \item a recording error,
    \item a legitimate rare event.
\end{itemize}

Therefore statistical evidence should trigger investigation,
not automatic conclusions.

\end{frame}

% ============================================================
% FINAL SUMMARY
% ============================================================

\section{Final Summary}

\begin{frame}{Complete Unit--IV Concept Map}
\small

\[
\boxed{
\text{WORK STUDY}
}
\]

\[
\downarrow
\]

\[
\boxed{
\text{METHOD STUDY}
}
\quad+\quad
\boxed{
\text{WORK MEASUREMENT}
}
\]

\vspace{5pt}

Method study:

\[
\text{Select}\rightarrow\text{Record}\rightarrow\text{Examine}
\rightarrow\text{Develop}\rightarrow\text{Install}\rightarrow\text{Maintain}
\]

\vspace{5pt}

Work measurement:

\[
OT\rightarrow NT\rightarrow ST
\]

\vspace{5pt}

Quality control:

\[
\text{Variation}
\rightarrow
\text{Control Charts}
\rightarrow
\text{Process Investigation}
\]

\end{frame}

% ------------------------------------------------------------

\begin{frame}{Unit--IV: Key Equations}
\small

\[
\boxed{
NT=OT\frac{R}{100}
}
\]

\[
\boxed{
ST=NT(1+A)
}
\]

\[
\boxed{
\hat p=\frac{x}{n}
}
\]

\[
\boxed{
n=\frac{z^2pq}{e^2}
}
\]

\[
\boxed{
UCL_{\bar X}=\bar{\bar X}+A_2\bar R
}
\]

\[
\boxed{
LCL_{\bar X}=\bar{\bar X}-A_2\bar R
}
\]

\[
\boxed{
UCL_R=D_4\bar R
}
\]

\[
\boxed{
UCL_p=\bar p+
3\sqrt{\frac{\bar p(1-\bar p)}{n}}
}
\]

\[
\boxed{
UCL_c=\bar c+3\sqrt{\bar c}
}
\]

\end{frame}

% ------------------------------------------------------------

\begin{frame}{Final Examination Checklist}
\small

Before an examination, be able to explain:

\begin{itemize}
    \item Work study and its objectives.
    \item Method study and SREDIM.
    \item ECRS.
    \item OPC, flow process, two-handed and man-machine charts.
    \item Work measurement.
    \item Stopwatch time study.
    \item Continuous and snapback timing.
    \item Observed time.
    \item Performance rating.
    \item Normal time.
    \item Allowances.
    \item Standard time.
    \item Work sampling.
    \item Binomial model.
    \item Confidence interval.
    \item Sample-size calculation.
    \item SQC.
    \item Common and special causes.
    \item Variables and attributes.
    \item $\bar X$--$R$ charts.
    \item $p$ and $c$ charts.
    \item Acceptance sampling.
    \item OC curve.
    \item AQL, LTPD, $\alpha$ and $\beta$.
\end{itemize}

\end{frame}

% ------------------------------------------------------------

\begin{frame}{Final Conceptual Summary}
\small

\begin{center}

\Large
\textbf{Improve the Method}\\[4pt]

\[
\downarrow
\]

\textbf{Measure the Work}\\[4pt]

\[
\downarrow
\]

\textbf{Establish Normal and Standard Time}\\[4pt]

\[
\downarrow
\]

\textbf{Monitor Process Variation}\\[4pt]

\[
\downarrow
\]

\textbf{Identify Statistical Signals}\\[4pt]

\[
\downarrow
\]

\textbf{Investigate and Improve the Process}

\end{center}

\vspace{5pt}

\[
\boxed{
\text{Productivity improvement and quality control are continuous engineering activities}
}
\]

\end{frame}

% ============================================================
% END
% ============================================================

\begin{frame}[standout]
\centering
\Huge
\textbf{Thank You}

\vspace{8pt}

\Large
Questions and Discussion
\end{frame}

\end{document}
